\chapterbackmatter{Solutions to Weinberg Exercises}
\ExerciseEditorialNote

\begin{enumerate}
  \WeinbergSolution{1}
  It is useful to do the calculation without choosing coordinates on the
  gauge group.  Put
  \[
    C=t_AC^A,\qquad V=t_AV^A,\qquad \Omega=t_A\Omega^A ,
  \]
  and count every Lie-algebra commutator as one power of a gauge coupling,
  as in Eq.~\eqref{eq:27.1.9}.  Restrict a superfield to the left-chiral
  slice and write
  \[
    V_L=C+\Delta,\qquad \Delta|_{\theta=0}=0 .
  \]
  A superfield in Wess--Zumino gauge vanishes on this slice.  The
  transformation law \eqref{eq:27.1.12} therefore gives
  \[
    e^{-i\Omega_L}=e^{-C}e^{2V_L},\qquad
    -i\Omega_L=\log\!\left(e^{-C}e^{2(C+\Delta)}\right),
    \tag{1}
  \]
  where we have made the convenient choice \(W=iC\).  A real part of
  \(W\) may be added; it is just the residual ordinary gauge
  transformation.

  The Baker--Campbell--Hausdorff formula through two commutators gives
  \[
    -i\Omega_L=C+2\Delta-[C,\Delta]
      +\frac12[C,[C,\Delta]]
      +\frac13[\Delta,[C,\Delta]]+O(g^3).
    \tag{2}
  \]
  Equation (2), together with the normalization of
  Eq.~\eqref{eq:27.1.14}, is already a compact expression for every
  required component.  To make the component prescription completely
  explicit, let \(\Pi_1\) and \(\Pi_2\) extract, respectively, the
  coefficients normalized as
  \[
    \Pi_1\Omega_L=w,\qquad \Pi_2\Omega_L=\mathcal W
  \]
  in Eq.~\eqref{eq:27.1.14}.  Then
  \[
  \begin{aligned}
    W&=iC,\\
    w&=\sqrt2\left\{-\omega+\frac12[C,\omega]
                 -\frac14[C,[C,\omega]]\right\}+O(g^3),\\
    \mathcal W
      &=\Pi_2\,i\left\{
          2\Delta-[C,\Delta]+\frac12[C,[C,\Delta]]
          +\frac13[\Delta,[C,\Delta]]\right\}+O(g^3).
  \end{aligned}
  \tag{3}
  \]
  The last line includes the bilinear-in-\(\omega\) contribution generated
  by the last BCH commutator; omitting it is the usual error in doing the
  calculation component by component.  If
  \(\Delta_1\) and \(\Delta_2\) denote the terms of degree one and two in
  the left-chiral Grassmann coordinate, the same result can be displayed
  without projection notation:
  \[
    (\Omega_L)_2=i\left\{
      2\Delta_2-[C,\Delta_2]+\frac12[C,[C,\Delta_2]]
      +\frac13[\Delta_1,[C,\Delta_1]]\right\}.
    \tag{4}
  \]
  In ordinary components put \(X=M-iN\).  The chiral-slice Fierz identity
  \[
    (\bar\theta P_L\omega_B)(\bar\theta P_L\omega_D)
      =-\frac12(\bar\theta P_L\theta)
        (\omega_{BL}^{\mathsf T}\epsilon\omega_{DL})
  \]
  turns (4) into
  \[
    \mathcal W=X-\frac12[C,X]+\frac14[C,[C,X]]
      +\frac i6[t_B,[C,t_D]]
        (\omega_{BL}^{\mathsf T}\epsilon\omega_{DL})+O(g^3).
    \tag{5}
  \]
  At zeroth order Eqs.~(3)--(5) reduce to
  \[
    W=iC,\qquad w=-\sqrt2\,\omega,\qquad
    \mathcal W=M-iN,
  \]
  as follows directly from Eq.~\eqref{eq:27.1.17}.  Complex conjugation
  supplies the right-chiral components.  Substitution of (2) back into
  (1) gives
  \(C'=\omega'=M'=N'=O(g^3)\), while \(V_\mu,\lambda,D\) remain as the
  Wess--Zumino-gauge fields.

  \WeinbergSolution{2}
  Begin with two otherwise independent chiral spinor superfields.  Lorentz
  covariance decomposes the coefficient linear in the chiral Grassmann
  coordinate into a scalar, a pseudoscalar, an antisymmetric tensor, and
  its dual.  Apply Eq.~\eqref{eq:27.2.20} successively at Grassmann degrees
  zero, one, and two.  The degree-zero equation identifies the two lowest
  spinors as the chiral projections of one Majorana spinor \(\lambda\).
  The scalar part of the degree-one equation leaves one real auxiliary
  scalar \(D\), while its tensor part identifies a single real
  antisymmetric tensor \(f_{\mu\nu}\).  The degree-two equation fixes the
  highest spinors as derivatives of \(\lambda\).

  Thus the most general solution is
  \[
  \begin{aligned}
    W_L={}&\lambda_L(x_+)
      +\frac12\gamma^\mu\gamma^\nu\theta_L f_{\mu\nu}(x_+)
      -i\theta_LD(x_+)
      +(\theta_L^{\mathsf T}\epsilon\theta_L)
         \sl{\partial}\lambda_R(x_+),\\
    W_R={}&\lambda_R(x_-)
      +\frac12\gamma^\mu\gamma^\nu\theta_R f_{\mu\nu}(x_-)
      -i\theta_RD(x_-)
      -(\theta_R^{\mathsf T}\epsilon\theta_R)
         \sl{\partial}\lambda_L(x_-).
  \end{aligned}
    \tag{6}
  \]
  Here \(D=D^*\), \(f_{\mu\nu}=f_{\mu\nu}^*\), and
  \(\lambda\) is Majorana.  In the general-spinor-superfield notation the
  remaining components are consequently
  \[
  \begin{gathered}
    C_{(\alpha)}=\lambda_\alpha,\qquad
    \omega_{(\alpha)\beta}
      =\frac i2(\gamma^\mu\gamma^\nu\epsilon)_{\alpha\beta}f_{\mu\nu}
       +(\gamma_5\epsilon)_{\alpha\beta}D,\\
    V_{(\alpha)\mu}=-i\partial_\mu(\gamma_5\lambda)_\alpha,\quad
    M_{(\alpha)}=-i(\sl{\partial}\gamma_5\lambda)_\alpha,\quad
    N_{(\alpha)}=-\sl{\partial}\lambda_\alpha,\quad
    \lambda_{(\alpha)\beta}=D_{(\alpha)}=0 .
  \end{gathered}
    \tag{7}
  \]

  The tensor part that is not accounted for by the real
  \(f_{\mu\nu}\) is precisely its exterior derivative.  Equation
  \eqref{eq:27.2.20} sets it to zero:
  \[
    \epsilon^{\mu\nu\rho\sigma}\partial_\rho f_{\mu\nu}=0.
    \tag{8}
  \]
  This is the homogeneous Maxwell equation.  Locally, on a contractible
  patch, the Poincare lemma then gives
  \(f_{\mu\nu}=\partial_\mu V_\nu-\partial_\nu V_\mu\).

  \WeinbergSolution{3}
  Label the adjoint chiral multiplet by
  \((\phi_A,\psi_A,\mathcal F_A)\), \(A=1,2,3\), and take
  \[
    (t_A)_{BC}=-ig\epsilon_{ABC},\qquad
    C_{ABC}=g\epsilon_{ABC}.
  \]
  There is no invariant vector, and the only invariant symmetric tensor of
  rank at most three is \(\delta_{AB}\).  The apparent cubic invariant
  \(\epsilon_{ABC}\Phi_A\Phi_B\Phi_C\) vanishes because scalar
  superfields commute.  Apart from an irrelevant constant, the most general
  renormalizable superpotential is therefore
  \[
    f(\Phi)=\frac m2\sum_A\Phi_A\Phi_A ,
    \tag{9}
  \]
  with complex \(m\).

  In the conventions of Eq.~\eqref{eq:27.4.1} the complete off-shell
  Lagrangian is
  \[
  \begin{aligned}
    \mathcal L={}&
      -(D_\mu\phi)_A^*(D^\mu\phi)_A
      -\frac12\bar\psi_A(\sl D\psi)_A
      +\mathcal F_A^*\mathcal F_A
      +2\operatorname{Re}(m\phi_A\mathcal F_A)
      -\operatorname{Re}\!\left[
        m(\psi_{AL}^{\mathsf T}\epsilon\psi_{AL})\right]\\
    &-2\sqrt2\operatorname{Re}\!\left[
       g\epsilon_{ABC}
       (\lambda_{AL}^{\mathsf T}\epsilon\psi_{CL})\phi_B^*\right]
      +\frac12D_AD_A
      +ig\epsilon_{ABC}\phi_B^*\phi_CD_A\\
    &-\frac14f_{A\mu\nu}f_A^{\mu\nu}
      -\frac12\bar\lambda_A(\sl D\lambda)_A
      +\frac{g^2\vartheta}{64\pi^2}
        \epsilon_{\mu\nu\rho\sigma}f_A^{\mu\nu}f_A^{\rho\sigma},
  \end{aligned}
    \tag{10}
  \]
  where repeated adjoint indices are summed and
  \[
  \begin{aligned}
    (D_\mu\phi)_A&=\partial_\mu\phi_A
       +g\epsilon_{ABC}V_{B\mu}\phi_C,\\
    (D_\mu\psi)_A&=\partial_\mu\psi_A
       +g\epsilon_{ABC}V_{B\mu}\psi_C,\\
    f_{A\mu\nu}&=\partial_\mu V_{A\nu}-\partial_\nu V_{A\mu}
       +g\epsilon_{ABC}V_{B\mu}V_{C\nu}.
  \end{aligned}
  \]
  The auxiliary equations are
  \[
    \mathcal F_A=-m^*\phi_A^*,\qquad
    D_A=-ig\epsilon_{ABC}\phi_B^*\phi_C .
    \tag{11}
  \]
  Hence
  \[
    V=|m|^2\phi_A^*\phi_A
      +\frac12\sum_A
       \left|-ig\epsilon_{ABC}\phi_B^*\phi_C\right|^2\geq0.
    \tag{12}
  \]
  The origin always has \(V=0\), so supersymmetry cannot be broken.

  For \(m\ne0\) the origin is the unique vacuum.  The \(SU(2)\) symmetry is
  unbroken: the three vector bosons and three gauginos are massless, while
  the three adjoint chiral multiplets have mass \(|m|\).  For \(m=0\),
  Eq.~(11) vanishes when the real and imaginary parts of \(\boldsymbol\phi\)
  are parallel:
  \[
    \boldsymbol\phi=e^{i\alpha}v\,\boldsymbol n,\qquad
    \boldsymbol n^2=1 .
    \tag{13}
  \]
  At \(v\ne0\), rotate \(\boldsymbol n\) to the third direction.
  Then \(SU(2)\to U(1)\); the two broken-direction vector multiplets have
  mass \(\sqrt2\,g|v|\), while the third vector multiplet and the neutral
  chiral modulus are massless.  At \(v=0\) all fields are massless.

  \WeinbergSolution{4}
  First keep the normalization independent of the particular parametrization
  of Eq.~\eqref{eq:27.5.1}.  Let the two charges be \(+q\) and \(-q\), let
  \(v\) be the nonzero scalar expectation value in the Higgs phase, and put
  \[
    a=\sqrt2\,|qv|,\qquad M_F=\sqrt{|m|^2+a^2}.
    \tag{14}
  \]
  Phases of \(m\) and \(v\) may be removed by field redefinitions; take
  them real below.

  In the unbroken-gauge phase, \(v=0\).  Both chiral auxiliaries vanish,
  while \(D_0\ne0\), so Eq.~\eqref{eq:27.5.12} says simply
  \[
    \lambda_L=G_L,\qquad
    \Psi_L=\psi_{-L},\qquad
    (\Psi^c)_L=\psi_{+L},
    \tag{15}
  \]
  up to an overall phase of the normalized goldstino \(G\).
  The two chiral spinors form a Dirac fermion \(\Psi\) of mass \(|m|\);
  the gaugino is the massless goldstino.

  For definiteness take the negatively charged scalar to condense,
  \(\phi_{-0}=v\), \(\phi_{+0}=0\).  The fermion mass terms couple
  \(\psi_{-L}\) to the combination
  \(m\psi_{+L}-ia\lambda_L\).  Define
  \[
    X_L=\frac{m\psi_{+L}-ia\lambda_L}{M_F},\qquad
    G_L=\frac{-ia\psi_{+L}+m\lambda_L}{M_F},\qquad
    Y_L=\psi_{-L}.
    \tag{16}
  \]
  Then \(G\) is massless, and \(X,Y\) are the two chiral projections of a
  Dirac fermion of mass \(M_F\).  Equivalently, the original fields are
  \[
    \psi_{+L}=\frac{ia\,G_L+mX_L}{M_F},\qquad
    \lambda_L=\frac{mG_L+iaX_L}{M_F},\qquad
    \psi_{-L}=Y_L .
    \tag{17}
  \]
  The relative phase in (16) agrees with the \(i\sqrt2\mathcal F\) and
  \(D\) direction in Eq.~\eqref{eq:27.5.12}; changing the phase convention
  for \(\lambda\) changes both displayed \(i\)'s.  For the opposite sign of
  \(\xi\), interchange \(+\leftrightarrow-\) and complex conjugate the
  charge phase.  Equations (14)--(17) exhibit explicitly that the
  goldstino is the unique zero mode in either phase.

  \WeinbergSolution{5}
  Regard the complex adjoint scalar as a traceless \(3\times3\) matrix
  \(\phi\).  With no hypermultiplets, Eq.~\eqref{eq:27.9.10} gives
  \[
    V=0\quad\Longleftrightarrow\quad
    [\operatorname{Re}\phi,\operatorname{Im}\phi]=0.
    \tag{18}
  \]
  The two Hermitian matrices may therefore be diagonalized simultaneously.
  Modulo an \(SU(3)\) gauge transformation and the Weyl group,
  \[
    \phi_\infty=\operatorname{diag}(a_1,a_2,a_3),\qquad
    a_1+a_2+a_3=0,
    \tag{19}
  \]
  where the two independent \(a_i\) are arbitrary complex numbers.

  If the three eigenvalues are distinct, the only massless gauge fields are
  the two Cartan photons, so \(SU(3)\to U(1)\times U(1)\).  If two
  eigenvalues coincide, the corresponding root gauge fields are also
  massless and the unbroken group is \(SU(2)\times U(1)\).  All eight gauge
  fields are massless only at \(a_1=a_2=a_3=0\).

  Choose orthonormal Cartan generators \(H_I\), \(I=1,2\), and write
  \(\phi_\infty=v_IH_I\).  Define the electric and magnetic fluxes of the
  two massless Cartan fields by
  \[
    q_I=\int_{S_\infty^2}dS_i\,E_I^i,\qquad
    \mathcal M_I=\int_{S_\infty^2}dS_i\,B_I^i .
    \tag{20}
  \]
  The surface-term calculation leading to
  Eq.~\eqref{eq:27.9.22} applies independently in the two commuting
  directions and gives
  \[
    Z_{12}=2\sqrt2\sum_{I=1}^2
       v_I\bigl(iq_I-\mathcal M_I\bigr).
    \tag{21}
  \]
  This form is basis independent: a simultaneous orthogonal change of the
  Cartan basis rotates \(v_I,q_I,\mathcal M_I\) and leaves their
  contractions unchanged.  On a degeneracy wall the same formula holds,
  with \(q_I,\mathcal M_I\) understood as the charges in the Abelian
  factors of the unbroken group.
\end{enumerate}
